% ----------
% Tech report template
% 
% ----------
\documentclass[10pt,a4paper]{article}
\usepackage{geometry}
%\usepackage[utf8x]{inputenc}
\usepackage[T1]{fontenc}
\usepackage[english]{babel}
\usepackage[runin]{abstract}
\usepackage{titling}
\usepackage{listings}
\usepackage{booktabs}
\usepackage{mlmodern}
\usepackage{enumitem}
%\usepackage{piton} % Without this enable, it might not work for the listing package. Switch to LuaLateX if you want to use this option, however. 
\usepackage{fancyhdr}
%\usepackage{helvet}
\usepackage{csquotes}
\usepackage{graphicx}
\usepackage{blindtext}
\usepackage{parskip}
\usepackage{etoolbox}
%\usepackage[scaled=0.90]{helvet} % Font 1
%\usepackage[scaled=0.85]{beramono} % Font 2
\usepackage{xcolor}
\usepackage{hyperref}
% Bibliography
\usepackage[backend=biber,style=alphabetic]{biblatex}
\addbibresource{bibb.bib} %Imports bibliography file

\input{preamble.tex}

% ----------
% Variables
% ----------

\title{\textbf{Research topic inquiry and analysis record},\\ v1, first edition} % Full title of your tech report
\runningtitle{Research Record Studies} % Short title
\author{Fujimiya Amane} % Full list of authors
\runningauthor{Amane et al.} % Short list of authors
\affiliation{} % Affiliation e.g. University or Company
\department{Your Department} % Department or Office
\memoid{RES-RP-XX} % ID of the tech report
\theyear{Year} % year of the tech report
\mydate{Full month, year} %the date
\usepackage{tcolorbox}
\usepackage{fancyvrb}

\tcbuselibrary{minted,breakable,xparse,skins}

\definecolor{bg}{gray}{0.95}
%\DeclareTCBListing{mintedbox}{O{}m!O{}}{%
%  breakable=true,
%  listing engine=minted,
%  listing only,
%  minted language=#2,
%  minted style=default,
%  minted options={%
%    linenos,
%    gobble=0,
%    breaklines=true,
%    breakafter=,,
%    fontsize=\small,
%    numbersep=8pt,
%    #1},
%  boxsep=0pt,
%  left skip=0pt,
%  right skip=0pt,
%  left=25pt,
%  right=0pt,
%  top=3pt,
%  bottom=3pt,
%  arc=5pt,
%  leftrule=0pt,
%  rightrule=0pt,
%  bottomrule=2pt,
%  toprule=2pt,
%  colback=bg,
%  colframe=orange!70,
%  enhanced,
%  overlay={%
%    \begin{tcbclipinterior}
%    \fill[orange!20!white] (frame.south west) rectangle ([xshift=20pt]frame.north west);
%    \end{tcbclipinterior}},
%  #3}

% Python preset environment. 
\newcommand{\pyobject}[1]{\fbox{\color{red}{\texttt{#1}}}}

\NewDocumentCommand{\codeword}{v}{%
\texttt{\textcolor{blue}{#1}}%
}
% Inline python highlighting (Might not look great, but eh)
\lstset{language=Python,keywordstyle={\bfseries \color{blue}}}


% ----------
% actual document
% ----------
\begin{document}
\maketitle

\begin{abstract}
    \noindent
    This is a template for all RHINELAB-RINALAB \textbf{research topic inquiry} and analysis writing. Please use this template and specify information regarding what is required. This is a template for \textbf{writing}, so it is less strict than the majority of all templates. 
    
    This is \textbf{version 1} of the template. Rather would there not be a lot of changes, but just so happens, then here. 

    \keywords{Keyword1, Keyword2, Keyword3}
\end{abstract}

\vspace{2.5cm}



\thispagestyle{firstpage}

\pagebreak

% ----------
% End of first page
% ----------

\newgeometry{} % Redefine geometries (normal margins)

\section{Specification}
\label{sec:spec}

Provide specification details, metaprogramming. 

\subsection{Document code, priority}

Provide document codes, additional document code, purpose, organization, authorship and anyone contributing to the report, identification of purpose (proposal, or research, indexing 01). Priority code, additional conditions. Does it override anything, or add to anything?

\subsection{Body of issue}

Who issue, team or departmental attachment. 

\subsection{Date and Version}

Version specification across versions, and similarly date of construction. \textbf{Do not remove old dates}, for versioning record. 

\section{Introduction}
\label{sec:intro}

The introduction is to present:

\begin{itemize}[topsep=1pt,itemsep=1pt]
    \item Statement of the topic in particular, or research topic. 
    \item Description of it. Field specification and whatever information regarding the \textbf{scope} of the topic analysis. 
\end{itemize}

Usually this is the only thing required. 

\section{Details}

Interjection and breakdowns, detailed analysis and so on of the topic in question. Mostly is to specify: 

\begin{itemize}
    \item The full chain of various \textbf{observations}, analysis, \textbf{thinking, logical arguments}, with \textbf{realization} of concepts of the (research) topic. 
    \item \textbf{Problems} either born out pre-analysis or during analysis. Statement, descriptions and so on of that research topic or discussion topic. 
    \item \textbf{Questions} regarding all aspects of the research analysis, including also \textit{question of problems} in regard to what exist. 
    \item Series of \textbf{resolution} about the topic itself. Note that it is not \textbf{required to be perfect and complete}, absolutely not so. You do not resolve the problem here. If you want to, this is already a paper elsewhere.  
\end{itemize}

Those would be many conditions required of this section. You can also leave observations and thus. Indeed, one can also leave some \textbf{prototype} and more details. While this part is a named section, you should not be constrained to express the report this way, but \textbf{change it, with condition that this 'idea of the section like this'} exists in the final document. 

If we are to say the principle here. 

\begin{quote}
    You \textbf{must} document the mess, the open questions, the partial resolutions. The template won't let you present only the clean story, and \textbf{never} will you go past this.
\end{quote}

\subsection{References and resources}

Provide all resources, references taken during analysis, and any form of results, preliminary probing, and statements about availability of demos and so on if there exists any.

\clearpage
% ----------
% Bibliography
% ----------

% \bibliography{../biblio.bib}
% \bibliographystyle{abbrv}

\appendix

\section{An appendix}

Appendix to more information. Analysis, cost process, etc. 

\section{Another appendix}

\subsection*{Subsectioning in appendix}

Some more text, and a list:

\blindenumerate



\section{Citations}
Items that are cited: \textit{The \LaTeX\ Companion} book \cite{latexcompanion} together with Einstein's journal paper \cite{einstein} and Dirac's book \cite{dirac}---which are physics-related items. Next, citing two of Knuth's books: \textit{Fundamental Algorithms} \cite{knuth-fa} and \textit{The Art of Computer Programming} \cite{knuthwebsite}.


Remember to change this to your respective bibliography. Usually, the file that holds the bibliography in the path is \texttt{bibb.bib}, so beware of that. The below part where there is \texttt{printbibliography} does it work - it prints the bibliography. Cite them using the \texttt{cite} key seen in the above citation. 

\medskip

\printbibliography %Prints bibliography


\end{document}

% ----------
